% P10-4b, 14.09.2026: Begriffs-/Beleganschlüsse und abgegrenzte Prüfaussagen präzisiert.
% M09, 11.09.2026: Leserführung und MCR konzentriert; neue
% HERKUNFT-M09-Kommentare unten sind die aktuelle Schreibgrundlage.
% Die folgende Versionshistorie beschreibt frühere Fassungen.
% ============================================================
% §5.4 Kontrollierte semantische Transformation — v05 (Claude, 04.09.)
% v05 (04.09., v23-Lesungs-Runde, P2): Analyst-Absatz VOR den
%   abschließenden Verdikt-/Aussagegrenzen-Block gezogen
%   (Dramaturgie: Muster -> Fall -> Analyst -> gemeinsame
%   Aussagegrenze -> Übergang; kein inhaltlicher Umbau).
% v04 (04.09., ANALYST-EINBAU + Schlussrunde): Core-Analyst-Absatz
%   am Abschnittsende (disk-analyst.md §4-①; Scope-Kipp per
%   Autor-Entscheid); Schlussrunden-Fixes „dieser
%   Verantwortungsverteilung" + „Analysefunde/konditionale
%   Autorisierung". Danach wieder SEMANTISCH FROZEN.
% v03 (Abschluss-Punkt): „Abdeckungsverstöße" → „Verstöße gegen eine
% strukturelle Zuordnungsregel zwischen Themenclustern und
% vorgeschlagenen Arbeitspaketen" — Kollisions-Schutz zum Ledger-
% Coverage-Begriff (5.3/Kap. 7); deckt sich mit der Kap.-4-Abgrenzung
% von E45-02 („struktureller Abdeckungs-Check, keine
% Quellvollständigkeit"). Kollegen-Urteil danach: §5.4 FROZEN
% (Titel-Doppelung 4.5/5.4 bewusst OK: Herleitung vs. Architektur);
% Freeze-Entscheid = Autor. Weiter mit 5.5 (Kollegen-Votum: KEIN
% Governance-Probeschnitt-Einschub).
% v02 (Kollegen-Review, alle vier Punkte übernommen): ① MCR EXPLIZIT
% benannt („im Folgenden als Maker--Checker--Repair (MCR) bezeichnete
% Quermuster") — der Abschnitt ETABLIERT den Begriff fürs EINMAL-
% Prinzip (5.6 verweist zurück!); „lebt" entfernt ② „Bestehendes
% angeglichen" → „anhand von Prüfbefunden überarbeitet" (klang nach
% 5.5-Fortschreibung) ③ Repariert wird der VORSCHLAG, nicht der
% Befund; „Das Modell darf Bedeutung erzeugen" → „Modellgestützte
% Komponenten übernehmen die semantische Erzeugung" (nüchtern, nicht
% anthropomorph) ④ E45-02-Grammatik („korrekt WAR oder … wirksam
% ist", eigener Satz) + „Gesamtqualität". E45-02-Referenzform: E-IDs
% sind als Belegschlüssel in Kap. 4 + Anhang-Ankertabelle formal
% eingeführt (Erklärsatz dort) — Arbeitsmarker bleibt korrekt.
% Kollegen-Urteil: danach einfrieren („Freeze-Stand 95 %") —
% Freeze-Entscheid = Autor.
% Grundlage: schema-mcr-schreibvertrag-5-4.md v02 (Wahrheitsliste
% #1–#10 + Verifikations-Protokoll §1b, code-geprüft 02.09.) ·
% Kollegen-Schreibauftrag 02.09. (vier Absätze, ~0,7–0,8 S.) ·
% finale Abbildung abb54musterformtransformation (gegen v02-Spec
% geprüft: Loop→Prüfung, H=Eskalation, Schranke=eigener Terminal-
% Marker, Triage NICHT in der Grafik).
%
% SELBSTKONTROLLE (5 Fragen des Auftrags, alle „nein"):
%   (a) keine Kap.-4-Genese (kein freier Prüfer/E45-03) (b) Muster
%   statt Instanzen (c) „bestanden" ≠ fachlich korrekt (Absatz 4)
%   (d) Repair-Zweig nur wo reparierbare Verdikte definiert (Absatz 3)
%   (e) Human Review und Versuchsschranke getrennt (Absatz 2).
% KEINE Klassen-/Enum-/Feld-/Promptnamen im Fließtext (GateLoop,
% GateDecision, Repairability etc. NUR hier in Kommentaren).
% Titel-Hinweis: Abschnittstitel ist identisch mit 4.5 (dort die
% HERLEITUNG von A4, hier die Säulen-Musterform) — Labels sind
% getrennt; falls der Autor Verwechslung fürchtet, wäre „Musterform
% der kontrollierten semantischen Transformation" die Alternative.
% Budget ~0,7–0,8 S. — 5.4 bleibt bewusst KLEINER als 5.3.
% ============================================================

\section{Kontrollierte semantische Transformation}
\label{sec:saeule-semantische-transformation}

Das \emph{Maker--Checker--Repair}-Muster (MCR) trennt die Erzeugung
eines fachlichen Vorschlags von seiner Prüfung und Überarbeitung.
Es setzt den entsprechenden Teil von A4 aus
Abschnitt~\ref{sec:kontrollierte-semantische-transformation} um und
wird innerhalb mehrerer Verarbeitungspfade wiederverwendet
(Abschnitt~\ref{sec:logische-gesamtarchitektur}). Es ist somit kein
zusätzlicher Schritt der Gesamtprozesskette. Menschliche
Autorisierung und Zustandsanwendung bleiben getrennte Aufgaben
(Abschnitt~\ref{sec:governance-mutation}).

\begin{figure}[htbp]
  \centering
  \includegraphics[width=0.9\textwidth]{figures/05/Abbildung5.4.pdf}
  \caption{Musterform der kontrollierten semantischen Transformation
    (eigene Darstellung).}
  \label{fig:musterform-transformation}
\end{figure}

Der \emph{Maker} erhält einen fachlichen Auftrag mit den für die
Aufgabe festgelegten Anforderungen und erzeugt einen Vorschlag.
Ein beendeter Modell- oder Agentenaufruf belegt noch nicht, dass
dieser Vorschlag die Abschlusskriterien der Stufe erfüllt.
Der \emph{Checker} prüft ihn deshalb anhand ausdrücklicher Kriterien:
deterministisch bei formal beschreibbaren Regeln, modellgestützt
bei inhaltlichen Urteilen. Das strukturierte Prüfergebnis -- in
Abbildung~\ref{fig:musterform-transformation} als \emph{Verdikt}
bezeichnet -- steuert den weiteren Workflow. Besteht der Vorschlag
die Prüfung, wird er weitergegeben. Liegt ein als reparierbar
eingestufter Befund vor und ist die Versuchszahl noch nicht
ausgeschöpft, erhält der modellgestützte \emph{Repair}-Schritt den
wörtlichen Prüfbefund als konkreten Überarbeitungsauftrag. Seine
neue Fassung durchläuft erneut dieselbe Prüfung. Die Rückmeldung
stammt somit aus einer gesonderten Prüfung; die Wiederholung allein
belegt noch keine Qualitätsverbesserung. Die Einstufung als
reparierbar erlaubt einen weiteren Versuch, garantiert aber nicht
dessen Erfolg.

Bei weiter bestehenden Fehlern endet die Schleife spätestens an
der festgelegten Versuchsschranke. Ohne einen reparierbaren Befund
wird keine weitere Modellüberarbeitung ausgelöst. Die betroffene
Stufe gibt dann einen dokumentierten Klärungs- oder Fehlerzustand
aus. Das Symbol H markiert den Bedarf an menschlicher Klärung;
ob dafür unmittelbar eine Anfrage entsteht oder die Stufe mit
einem Hinweis endet, legt der jeweilige Verarbeitungspfad fest.
Diese Klärung ist von der Freigabe eines bereits geprüften
Änderungsvorschlags zu unterscheiden. Jeder Prüfversuch der
Schleife wird protokolliert.
% HERKUNFT M09: GateLoop.Decide (Pass vor Schranke vor Repair),
% CanonicalCheckExecutor und IngestionHitlFinalizeExecutor:
% HumanReview/MaxAttemptsReached bedeuten dort terminale Ausgabe,
% keinen RequestPort-Aufruf. HumanReview ist kein Freigabeereignis.

Die gemeinsame Entscheidungsform wird unter anderem bei der
Themenbündelung, dem Zuschnitt von Arbeitspaketen und der
Kanonisierung des Ledgers verwendet. Ein Reparaturzweig ist nur
dort wirksam, wo die jeweilige Prüfung reparierbare Befunde
definiert. Ein konkreter Backlog-Fall zeigt das Zusammenspiel
(E45-02). Jede als Kern eines Themenclusters markierte Anforderung
sollte mindestens einem Arbeitspaket zugeordnet sein. Der Agent
konnte Entwürfe über ein Prüfwerkzeug kontrollieren und erhielt
dessen Befunde unmittelbar als Rückmeldung. Nach Ende seiner
Bearbeitung prüfte zusätzlich der Workflow das eingereichte Ergebnis.

Im ersten Versuch endete der Agentenaufruf, ohne einen Gesamtentwurf
zu speichern. Am Workflow-Prüfpunkt fehlten damit alle siebzig
geforderten Zuordnungen. Diese Befunde wurden dem erneut aufgerufenen
Agenten als Reparaturauftrag übergeben. Innerhalb dieses zweiten Versuchs prüfte der Agent einen Entwurf
mit 19 Arbeitspaketen. Das Werkzeug meldete darin noch vier
fehlende Verbindungen. Die anschließend gespeicherte Fassung
enthielt dieselben 19 Arbeitspakete mit den vier ergänzten
Verbindungen und bestand die erneute Workflow-Prüfung.
Die erhaltenen Werkzeugargumente zeigen zwischen diesen beiden
Fassungen unveränderte Titel, Ziele, Umfangsbeschreibungen und
Akzeptanzkriterien. Belegt sind damit die Rückkopplung, die
Ergänzung der Verbindungen und die Erfüllung der Strukturregel.
Ob die neuen Zuordnungen fachlich richtig sind, wurde nicht
semantisch geprüft; ebenso wenig ist dadurch die Vollständigkeit
gegenüber der ursprünglichen Quelle nachgewiesen.
% HERKUNFT M09: Lauf 20260804_074419_acab74,
% backlog-attempts.json: 70 UNCOVERED_CORE -> Repair -> Pass.
% ReClarifyBacklogGate.Check definiert die Zuordnungsregel.
% P10-4b: logs/otel-traces.jsonl Z.181 enthält den historischen
% Reparaturauftrag mit 70 wörtlichen Fehlermeldungen; Z.225/229
% enthalten den 19-PBI-Entwurf vor/nach vier ergänzten Verbindungen.
% Die frühere Annahme eines nicht archivierten Auftrags ist korrigiert.

Die Trennung von Interpretation und fachlicher Geltung lässt auch
offenere Analysen zu. Der Core-Analyst untersucht den Projektzustand
aus mehreren Perspektiven auf mögliche Lücken und Widersprüche.
Seine Funde werden mit Herkunft und Herleitung als strukturierter
Eingang für die reguläre Fortschreibung bereitgestellt. Erst dort
werden sie gegen den Bestand eingeordnet und gegebenenfalls zur
Autorisierung einer Änderung vorgelegt. Der Analyseworkflow selbst
mutiert den Core nicht; er ist keine weitere Instanz der hier
gezeigten Reparaturschleife.
% HERKUNFT M09: AnalystWorkflow/AnalystDeltaBuilder; eigener
% Fan-out/Fan-in-Graph mit Kritiker, kein GateLoop-Reparaturzweig.

Ein bestandenes Prüfergebnis belegt nur die Erfüllung der
festgelegten Kriterien. Kapitel~\ref{chap:evaluation} trennt daher
Kontrollnachweise, tatsächlich ausgelöste Reparaturschritte und
inhaltliche Bewertung. Die technische Komposition zeigt
Abschnitt~\ref{sec:realisierung-agent-executor-workflow}; der
folgende Abschnitt erklärt den Projektzustand, auf den sich die
fachlichen Vorschläge beziehen.
